\chapter{Epidermolysis Bullosa}
\index{Epidermolysis Bullosa}
\label{sec:Epidermolysis Bullosa}

\section{Overview}
Epidermolysis bullosa (EB) is a group of inherited blistering disorders characterized by mechanical fragility of the skin and mucous membranes, with an overall prevalence of 1 in 50,000. This condition affects a significant number of people worldwide and can range from mild to severe in presentation. Prompt diagnosis and appropriate management are essential for optimal outcomes. Early recognition and appropriate management are essential for optimal outcomes. A thorough understanding of the underlying mechanisms guides treatment decisions. Patient education and self-management play important roles in long-term care.
\section{Causes}
The underlying mechanisms involve complex interactions between genetic predisposition and environmental factors. Disruption of normal physiological processes leads to the characteristic features of the condition. Multiple contributing factors may act synergistically to trigger or worsen the condition. Understanding the specific cause helps guide targeted treatment approaches.
\section{Risk Factors}
Genetic predisposition and family history Advancing age Lifestyle factors (diet, exercise, smoking, alcohol) Environmental exposures Underlying medical conditions Medication use Stress and psychological factors Socioeconomic factors

\begin{itemize}[noitemsep]
  \item Genetic predisposition and family history
  \item Advancing age
  \item Lifestyle factors (diet, exercise, smoking, alcohol)
  \item Environmental exposures
  \item Underlying medical conditions
  \item Medication use
\end{itemize}
\section{Signs and Symptoms}
You may experience blister formation at sites of minor friction or trauma, ranging from mild acral blistering (EBS) to widespread blistering, scarring, milia, nail dystrophy, and mitten-hand deformities from repeated scarring (severe DEB). Extracutaneous involvement includes oral, esophageal, and laryngeal blistering, corneal erosions, anemia, growth retardation, and increased risk of aggressive squamous cell carcinoma in severe DEB. Symptoms vary depending on the severity and extent of the condition Pain or discomfort in the affected area Functional impairment affecting daily activities Changes in appearance or function of affected tissues Systemic symptoms may include fatigue, fever, or weight changes Symptoms may develop gradually or appear suddenly

\begin{itemize}[noitemsep]
  \item Pain or discomfort in the affected area
  \item Functional impairment affecting daily activities
  \item Changes in appearance or function of affected tissues
  \item Systemic symptoms may include fatigue, fever, or weight changes
  \item Symptoms may develop gradually or appear suddenly
\end{itemize}
\section{Progression and Stages}
This condition happens because of mutations in genes encoding structural proteins of the skin, classified by the level of blister formation within the basement membrane zone: EB simplex (EBS, keratin 5/14 mutations, intraepidermal cleavage), junctional EB (JEB, laminin 332/collagen XVII mutations, lamina lucida cleavage), and dystrophic EB (DEB, collagen VII mutations, sublamina densa cleavage). The condition may progress through different stages of severity over time. Early stages often involve mild or subtle symptoms that may be overlooked. Moderate stages involve more noticeable symptoms affecting daily function. Advanced stages may involve significant impairment and complications.
\section{Complications}
\begin{itemize}[noitemsep]
  \item Disease progression leading to worsening symptoms
  \item Secondary infections or injuries
  \item Side effects of treatment
  \item Reduced quality of life
  \item Emotional and psychological impact
\end{itemize}
\section{Diagnosis}
Doctors diagnose this based on clinical features, skin biopsy for immunofluorescence mapping, transmission electron microscopy, and genetic testing. Comprehensive medical history and physical examination Laboratory tests to assess organ function and rule out other conditions Imaging studies to visualize affected structures Specialized testing depending on the affected system Consultation with relevant specialists Monitoring of disease progression over time

\begin{itemize}[noitemsep]
  \item Comprehensive medical history and physical examination
  \item Laboratory tests to assess organ function
  \item Imaging studies to visualize affected structures
  \item Specialized testing depending on the affected system
  \item Consultation with relevant specialists
\end{itemize}
\section{Prevention}
Treatment typically involves multidisciplinary: gentle skin care, wound dressings (non-adherent), prevention and treatment of infection, nutritional support, esophageal dilation for strictures, physical therapy for contractures, and surveillance for squamous cell carcinoma.

\begin{itemize}[noitemsep]
  \item Maintain a healthy lifestyle with balanced diet and exercise
  \item Avoid known risk factors
  \item Attend regular medical check-ups for early detection
  \item Follow recommended screening guidelines
\end{itemize}
\section{Treatment Options}
Emerging therapies include gene therapy (beremagene geperpavec for DEB), protein replacement therapy, and cell-based therapies. Medications to manage symptoms and address underlying mechanisms Lifestyle modifications and self-care measures Physical therapy and rehabilitation Surgical intervention when indicated Regular monitoring and follow-up care Patient education and support services

\begin{itemize}[noitemsep]
  \item Medications to manage symptoms and address underlying mechanisms
  \item Lifestyle modifications and self-care measures
  \item Physical therapy and rehabilitation
  \item Surgical intervention when indicated
  \item Regular monitoring and follow-up care
\end{itemize}
\section{Living with It}
\begin{itemize}[noitemsep]
  \item Follow your treatment plan as prescribed
  \item Maintain regular follow-up with your healthcare provider
  \item Make necessary lifestyle adjustments
  \item Seek support from family, friends, or support groups
  \item Stay informed about your condition
\end{itemize}
\section{Outlook}
\begin{itemize}[noitemsep]
  \item Recovery varies widely: EBS has normal life expectancy
  \item Severe JEB and recessive DEB cause significant morbidity and mortality from infection, malnutrition, and aggressive squamous cell carcinoma.
\end{itemize}
